\chapter{자원: 회로는 공간을 차지한다}\label{ch:resources}

릴레이식 봄베에서 회로의 크기는 눈에 보인다. 6유닛에 릴레이가 3,455개 들고(테스트 레지스터 제외), 시프터 보드
한 장은 릴레이 130개를 싣느라 $250 \times 228$~mm가 되었으며, 기계의 폭은 1.88~m다. 켜져 있는 코일은 내내
전력을 열로 바꾸고(\ref{sec:el-power}절), 6유닛의 코일이 평균 0.21--0.31~kW를 쓴다. 회로를 하나 더하면 릴레이를 사서 기판에 꽂아야 하고,
기판과 전원과 팬이 함께 커진다. FPGA에서도 회로는 공간을 차지한다. 다만 그 공간은 칩 안에 미리 깔려 있는
부품의 수이고, 눈에 보이지 않는다.

회로의 비용은 둘이다. 부품이 차지하는 \emph{자원}과 신호가 지나는 데 걸리는 \emph{시간}. 이 장은 자원을 본다.
릴레이식에서 무엇이 릴레이를 많이 먹었는지를 먼저 떠올리고, FPGA 안에 어떤 부품이 있으며 코드가 그 부품을
얼마나 쓰는지 가늠하는 법으로 넘어간다. 이 감각이 없으면 \ref{ch:lesson}장의 사고가 난다.

\section{FPGA 안에 있는 것}

\Relay 릴레이식 봄베는 모든 보드를 극이 둘인 전환 접점 릴레이(DPDT) 한 종류로 짓는다(\ref{sec:relay-one}절).
논리는 접점의 직렬과 병렬로 만들고(\ref{sec:relay-logic}절), 1비트 기억은 자기 유지로 만든다(\ref{sec:relay-latch}절).
그래서 크기를 재는 단위가 릴레이의 개수다. 계수기 한 자리(5비트)가 릴레이 11개, 위치 표시 한 자리가 13개, 시프터
보드 한 장이 130개다(\ref{ch:relay-machine}장). 부품은 설계가 요구하는 만큼 사서 꽂는다. 그 대신 그것을 담을
그릇이 있다. 기계의 프레임에는 유닛 자리가 12개 있고, 먼저 만드는 6유닛은 그 절반을 채운다.

\Why FPGA는 만들어진 뒤에 회로를 바꿀 수 있는 칩이다. 게이트를 그때그때 새로 만들 수는 없으므로,
FPGA는 ``어떤 논리 함수로든 설정할 수 있는 작은 부품''을 수만 개에서 수십만 개 깔아 두고 그 사이를
설정 가능한 배선으로 잇는다. 그 작은 부품이 \term{LUT}{look-up table, 조회표}다. 릴레이식이 한 종류의 릴레이로
모든 것을 지은 것처럼, FPGA는 몇 종류 안 되는 범용 부품으로 모든 것을 짓는다. 다른 점은 부품의 수다. 칩마다
정해져 있고, 설계가 그 안에 들어가거나 들어가지 않거나 둘 중 하나다.

LUT는 이름 그대로 표다. 입력이 6개인 LUT(LUT6)는 $2^6 = 64$칸짜리 1비트 메모리이고, 입력 6비트를 주소로
삼아 그 칸의 값을 내놓는다. 64칸에 무엇을 적어 두느냐에 따라 입력 6개의 \emph{어떤} 논리 함수든 된다
(그림~\ref{fig:g-lut}). 합성 도구가 하는 일은 설계자가 적은 논리를 6입력 이하의 조각으로 쪼개 LUT에
하나씩 담는 것이다.

셈법이 릴레이와 다르다. 접점 회로에서는 식이 복잡할수록 접점이 늘었다. AND 하나에 직렬 접점 하나, OR 하나에
병렬 갈래 하나다. LUT에서는 식의 모양이 값을 정하지 않는다. 64칸 표에는 무엇이든 적을 수 있기 때문이다.

\begin{figure}[ht]
\centering
\begin{tikzpicture}[x=1cm,y=1cm]
  % 진리표
  \node[font=\small\bfseries] at (1.5,3.45) {LUT에 적힌 표 (일부)};
  \matrix[matrix of nodes,nodes={font=\ttfamily\footnotesize,minimum width=5mm,minimum height=4.2mm,inner sep=0pt},
          column sep=0pt,row sep=0pt] (m) at (1.5,0.9) {
     a & b & c & |[text=cgray]|\dots & {} & y \\
     0 & 0 & 0 & {} & {} & 0 \\
     0 & 0 & 1 & {} & {} & 1 \\
     0 & 1 & 0 & {} & {} & 0 \\
     0 & 1 & 1 & {} & {} & 1 \\
     1 & 0 & 0 & {} & {} & 0 \\
     1 & 0 & 1 & {} & {} & 1 \\
     1 & 1 & 0 & {} & {} & 1 \\
     1 & 1 & 1 & {} & {} & 1 \\
  };
  \draw (m-1-1.south west) -- (m-1-6.south east);
  \draw (m-1-5.north east) -- (m-9-5.south east);
  \node[note] at (1.5,-1.85) {\code{y = (a \& b) | c}};
  % LUT 기호
  \begin{scope}[xshift=6.6cm]
  \node[font=\small\bfseries] at (1.6,2.9) {LUT6 하나};
  \node[draw,fill=cbluebg,draw=cblue,minimum width=22mm,minimum height=30mm,align=center] (lut) at (1.6,0.8) {64칸\\메모리};
  \foreach \i/\l in {0/a,1/b,2/c,3/{},4/{},5/{}} {
    \draw[<-] ([yshift=11mm-\i*4.4mm]lut.west) -- ++(-0.8,0) node[left,font=\ttfamily\footnotesize] {\l};
  }
  \draw[->] (lut.east) -- ++(0.8,0) node[right,font=\ttfamily\footnotesize] {y};
  \node[note] at (1.6,-1.4) {입력 6비트가 주소, 출력은 그 칸의 값};
  \end{scope}
  % FF
  \begin{scope}[xshift=11.6cm]
  \node[font=\small\bfseries] at (1.2,2.9) {짝을 이루는 플립플롭};
  \node[reg,minimum width=12mm,minimum height=14mm] (ff) at (1.2,0.8) {FF};
  \pic at (ff.south) {clkmark};
  \draw[<-] (ff.west) -- ++(-0.6,0);
  \draw[->] (ff.east) -- ++(0.6,0);
  \node[note] at (1.2,-1.4) {1비트 기억};
  \end{scope}
\end{tikzpicture}
\caption{LUT는 작은 메모리로 만든 만능 게이트다. 그림~\ref{fig:g-assign}의 게이트 두 개는 LUT 하나에
들어간다(입력 세 개만 쓴다).}
\label{fig:g-lut}
\end{figure}

\General \code{y = (a \& b) | c}는 입력이 셋이므로 LUT 하나다. 접점으로 만들면 \code{a}와 \code{b}의 직렬에
\code{c}를 병렬로 단 접점 셋이지만, LUT에서는 게이트가 몇 개인지가 중요하지 않다.
중요한 것은 \emph{출력 하나가 입력 몇 개에 의존하는가}다. 입력이 여섯 이하이면 아무리 복잡한 식도
LUT 하나이고, 일곱 이상이면 LUT 여러 개를 엮어야 한다.

FPGA에는 LUT 말고도 몇 가지 부품이 있다(표~\ref{tab:g-res}).

\begin{table}[ht]
\centering
\caption{FPGA의 주요 자원. 수량은 이 프로젝트의 대상 칩 두 종의 대략적인 값이다.}
\label{tab:g-res}
\small
\begin{tabularx}{\linewidth}{l>{\raggedright\arraybackslash}X>{\raggedright\arraybackslash}p{27mm}>{\raggedright\arraybackslash}p{24mm}}
\toprule
자원 & 하는 일 & XC7A200T (Nexys Video) & XC7Z020 (PYNQ) \\
\midrule
LUT & 6입력 이하의 임의의 논리 함수. 조합 논리가 여기에 담긴다 & 약 13만 4천 & 약 5만 3천 \\
플립플롭(FF) & 1비트 기억. 레지스터가 여기에 담긴다 & LUT마다 딸려 있다 & \\
분산 RAM (LUT RAM) & LUT의 64칸 메모리를 그대로 작은 RAM으로 쓰는 것. 작고 빠른 표에 알맞다 & LUT의 일부를 전용 & \\
블록 RAM & 큰 메모리 전용 블록(수만 비트 단위). 읽는 데 한 클럭이 든다 & 365개 & 140개 \\
DSP 블록 & 곱셈기. 신호 처리에 쓰인다 & (이 프로젝트에서는 쓰지 않는다) & \\
\bottomrule
\end{tabularx}
\end{table}

\Ours 봄베의 회로에는 곱셈이 없고 큰 메모리도 필요 없다. 자원의 거의 전부가 LUT(전파 논리와 치환 계산)와
플립플롭(676비트 상태와 치환표)이다. 그래서 이 프로젝트에서 ``회로가 큰가''는 곧 ``LUT를 몇 개 쓰는가''다.

두 기계가 무엇을 세는지 나란히 놓으면 표~\ref{tab:g-res-count}와 같다. 릴레이 하나와 LUT 하나는 같은 크기의
부품이 아니므로 수를 직접 견줄 수는 없다. 볼 것은 \emph{어디에} 부품이 드는가다. 눈에 띄는 차이가 기억이다.
릴레이식은 676가닥의 상태를 기억하지 않는다. 전류가 한순간에 퍼지므로, 선의 상태 가운데 붙들어 두는 것은
결과를 읽는 테스트 레지스터의 26비트뿐이다(\ref{sec:t-register}절). FPGA는 한 걸음씩 퍼뜨리므로(\ref{sec:f-rounds}절) 676가닥 전부의 상태를
플립플롭에 담아 두어야 한다.

\begin{table}[ht]
\centering
\caption{두 기계의 크기를 재는 단위. 릴레이식은 6유닛의 값(보고서의 릴레이 수 표, 테스트 레지스터 제외), FPGA는
엔진 한 대를 합성한 단계의 값이다(\ref{ch:lesson}장).}
\label{tab:g-res-count}
\small
\begin{tabularx}{\linewidth}{l >{\raggedright\arraybackslash}X >{\raggedright\arraybackslash}X}
\toprule
& 릴레이식 & FPGA 엔진 \\
\midrule
부품 & 릴레이 3,455개. 그 가운데 시프터가 3,120개 & LUT 41,040개 \\
기억 & 자기 유지 릴레이. 선의 상태로 붙드는 것은 테스트 레지스터의 26비트 & 플립플롭 7,338개. 676가닥의 상태 전부를 포함 \\
표 & 배선 카드의 구리. 능동 부품이 없다 & LUT RAM 2,858개 \\
그릇 & 유닛 12자리의 프레임(폭 1.88~m) & 칩의 LUT 약 13만 4천 개 \\
따라오는 대가 & 코일의 전력과 열 (6유닛 평균 0.21--0.31~kW) & 칩이 찰수록 배선이 길어진다 (\ref{ch:pipeline}장) \\
\bottomrule
\end{tabularx}
\end{table}

\section{선택 회로는 비싸다}\label{sec:mux-cost}

회로에서 LUT를 가장 많이 잡아먹는 것은 덧셈도 비교도 아니고 \term{선택 회로}{multiplexer, mux}다.
여러 입력 가운데 하나를 골라 내보내는 회로다. 릴레이식에서도 사정이 같았다.

\Relay 릴레이식에서 ``여럿 가운데 하나를 고르는'' 회로는 두 번 나왔다. 하나는 위치 표시의 접점
나무다(\ref{sec:relay-tree}절). 5비트 값을 보고 램프 26개 가운데 하나에 전압을 보낸다. 전환 접점 하나가 길을
둘로 가르므로 26갈래에는 접점 25개가 든다. 고를 대상이 $n$개이면 접점이 $n-1$개, 대상의 수에 비례한다. 다른
하나는 이진 시프터다(\ref{sec:relay-shifter}절). 선마다 ``0--25칸 가운데 몇 칸 밀린 선으로 이을까''를 고르는 회로다.
위치마다 따로 배선하면 유닛 하나에 접점이 4,056개 들지만, 밀기를 1·2·4·8·16칸의 단으로 나누고 이웃한 밀기를
합쳐 1,040개(릴레이 520개)로 줄였다. 그렇게 줄이고도 6유닛의 릴레이 3,455개 가운데 3,120개가 시프터다.
릴레이식에서 가장 비싼 회로가 선택 회로였다.

\Why 선택 회로의 출력은 \emph{모든} 데이터 입력과 선택 신호에 의존한다. 넷 중 하나를 고르는 4:1 선택
회로는 데이터 4비트와 선택 2비트, 모두 여섯 입력이므로 LUT6 하나에 꼭 맞는다. 그보다 큰 선택 회로는
4:1을 층층이 쌓아야 한다. 16:1은 4:1 네 개로 16개를 4개로 줄이고 다시 4:1 하나로 고르므로 LUT 다섯 개다.
이런 식으로 $N$개 가운데 하나를 고르려면 LUT가 대략 $N/3$개 든다($N/4 + N/16 + \cdots$). 접점 나무와 같은
셈이다. 접점 나무의 갈림 하나는 2:1 선택이어서 $n-1$개가 들고, LUT 하나는 4:1 선택이어서 그 삼분의 일쯤 든다.
어느 쪽이든 고를 대상의 수에 비례한다. 그리고 이것은 \emph{출력 1비트당}의 값이다(그림~\ref{fig:g-muxtree}).
시프터의 한 단이 코일 신호 하나를 나눠 쓰면서도 26선마다 접점을 하나씩 따로 두는 것과 같다. 고르는 신호는
함께 쓰지만, 골라지는 데이터는 비트마다 제 길이 있어야 한다.

\begin{figure}[ht]
\centering
\begin{tikzpicture}[x=1cm,y=1cm]
  \node[font=\small\bfseries] at (0.9,2.6) {4:1 = LUT 1개};
  \node[mux,minimum width=20mm] (m0) at (0.9,0.9) {4:1};
  \foreach \i in {0,1,2,3} { \draw[<-] ([yshift=7.5mm-\i*5mm]m0.west) -- ++(-0.6,0); }
  \draw[->] (m0.east) -- ++(0.5,0);
  \draw[<-] (m0.south) -- ++(0,-0.45) node[below,font=\footnotesize] {선택 2비트};
  \begin{scope}[xshift=5.2cm]
  \node[font=\small\bfseries] at (2.0,2.6) {16:1 = LUT 5개};
  \foreach \k/\yy in {0/2.0,1/1.0,2/0.0,3/-1.0} {
    \node[mux,minimum width=9mm,minimum height=5mm,font=\scriptsize] (a\k) at (0.8,\yy) {4:1};
    \foreach \i in {0,1,2,3} { \draw ([yshift=3mm-\i*2mm]a\k.west) -- ++(-0.35,0); }
  }
  \node[mux,minimum width=20mm] (b) at (3.4,0.5) {4:1};
  \foreach \k in {0,1,2,3} { \draw[->] (a\k.east) -- ++(0.5,0) |- ([yshift=7.5mm-\k*5mm]b.west); }
  \draw[->] (b.east) -- ++(0.5,0);
  \end{scope}
  \begin{scope}[xshift=11.3cm]
  \node[font=\small\bfseries] at (1.6,2.6) {크기의 어림 (출력 1비트당)};
  \node[font=\small,align=left,anchor=north west] at (-0.4,2.2) {
     \begin{tabular}{@{}rl@{}}
       4:1 & LUT 1개 \\
       16:1 & 5개 \\
       26:1 & 열 개 안팎 \\
       676:1 & 이백 개 이상 \\
     \end{tabular}};
  \end{scope}
\end{tikzpicture}
\caption{선택 회로의 크기. 고를 대상이 많아지면 그에 비례해 커진다. 접점 나무가 26글자에 접점 25개를 쓰는
것과 같은 셈법이다.}
\label{fig:g-muxtree}
\end{figure}

\paragraph{한 방향의 선택.} 다른 점이 하나 있다. 접점의 선택은 양방향이다. 시프터 보드 한 장이 가는 길의
밀기와 오는 길의 되돌림을 같은 릴레이의 두 극으로 함께 맡고, 같은 보드를 거꾸로 끼우면 반대로 민다
(\ref{sec:relay-shifter}절). FPGA의 선택 회로에는 방향이 있다(\ref{sec:hdl-direction}절). 입력 여럿 가운데 하나를
출력으로 내보낼 뿐, 출력 쪽에서 들어온 신호를 입력 쪽으로 돌려보내지 못한다. 거꾸로 지나는 길이 필요하면
선택 회로를 한 벌 더 둔다.

\General 선택 회로는 코드에서 눈에 잘 띄지 않는다. 다음 두 줄은 모두 선택 회로다. 특히 둘째 줄, 변수로
배열의 원소를 고르는 것은 프로그램에서는 가장 값싼 연산이지만 회로에서는 ``배열의 모든 원소를 선택 회로의
입력에 연결한다''는 뜻이다.

\begin{lstlisting}[language=SV,numbers=none]
assign y = sel ? a : b;      // 2:1 mux
assign z = table[index];     // (size of table):1 mux, for every bit of z
\end{lstlisting}

\Ours 봄베의 회로는 선택 회로투성이다. ``26글자 가운데 하나''를 고르는 일이 곳곳에 있기 때문이다.
로터를 통과하는 것은 배선표 26칸에서 한 칸을 고르는 일이고(5비트 폭의 26:1), 전파 논리에서 ``케이블 $a$의
$y$번 선''을 읽는 것도 선택이다. 코드~\ref{lst:g-spread}의 \code{live\_q[cable\_a + y]}가 그것인데,
여기서 고르는 대상은 26개가 아니라 676개다. 이 한 줄이 이 프로젝트의 가장 큰 사고의 원인이었다.
\ref{ch:lesson}장에서 다룬다.

릴레이식에는 이 선택이 없다. ``메뉴의 이 선은 케이블 $a$와 케이블 $b$를 잇는다''는 관람객이 메뉴 케이블을
잭에 꽂아 정한다(\ref{sec:bombe-one-page}절, \ref{sec:c-scrambler}절). 꽂힌 케이블은 구동 내내 그대로이므로 접점도 릴레이도 들지 않는다.
FPGA에는 꽂을 케이블이 없다. 메뉴를 설정값으로 받고, 그 설정값으로 케이블을 고르는 선택 회로가 생긴다.
사람이 손으로 하던 고르기가 회로가 되는 것이다. \ref{ch:lesson}장의 해결책은 이 고르기를 ``676개 가운데
하나''가 아니라 ``케이블 26개 가운데 하나, 그 안에서 선 26개 가운데 하나''로 나누는 것이다. 위치마다 따로 잇던
26가지를 1·2·4·8·16칸 다섯 단으로 나눈 시프터처럼, 큰 선택을 구조에 맞춰 작은 선택으로 나누면 싸진다.

\section{표는 메모리에 담는다}

\Relay 릴레이식에서 로터의 배선표는 배선 카드 $W$의 구리다(\ref{sec:relay-shifter}절,
\code{hardware/relay/wiring-cards}). 카드에는 커넥터와 장착 구멍 말고 부품이 없다. 로터 카드 한 장에 가는 길
26가닥과 오는 길 26가닥이 동선으로 그려져 있을 뿐이다. ``$i$번 핀으로 들어온 전류는 $W[i]$번 핀으로 나간다''가
그대로 표 조회다. 릴레이도 코일 전력도 들지 않고, 표를 바꾸려면 카드를 바꿔 꽂는다. 기어식에서는 같은 표가
드럼 기판의 안쪽 층에 있다(\ref{sec:gm-drum}절). 드럼으로 돌아오는 전류는 같은 배선을 거꾸로 지나므로
역방향 표 $W^{-1}$이 저절로 된다. 다만 두 기계 모두 \emph{계산해 둔} 표는 없다. 위치마다의 치환은 시프터의
접점이나 드럼의 회전이 그 자리에서 만들 뿐, 어디에도 적어 두지 않는다.

\Why 같은 계산을 되풀이하는 것보다 계산해 둔 결과를 표에서 읽는 편이 싸다. 프로그램에서도 쓰는 방법이지만
회로에서는 의미가 더 크다. 계산 회로는 쓰든 안 쓰든 자리를 차지하고 있는 반면, 메모리는 같은 양의 정보를
훨씬 좁은 자리에 담는다. FPGA에는 메모리가 두 종류 있다. LUT를 메모리로 쓰는 \term{분산 RAM}{distributed
RAM}은 작지만 주소를 주면 같은 클럭 안에 값이 나오고, \term{블록 RAM}{block RAM}은 크지만 읽는 데 한 클럭이
걸린다. FPGA의 표는 구리처럼 공짜가 아니다. 표 하나가 LUT나 플립플롭 몇 개를 차지한다. 그리고 거꾸로 흐르는
전류가 없으므로 역방향으로 읽을 표도 따로 담아야 한다.

\General 소리를 내는 회로가 사인파를 필요로 한다고 하자. 사인 함수를 계산하는 회로를 만드는 대신 한 주기의
값을 256칸짜리 표에 넣어 두고 차례로 읽어 낸다. 표는 블록 RAM 하나의 일부에 들어간다.

\Ours 로터의 배선표는 바뀌지 않는 상수 표다. 기준 RTL에서는 \code{wirings\_pkg}에 상수로 들어 있고
로터를 통과하는 함수가 그 표를 읽는다. FPGA 엔진은 구동을 시작할 때 고른 로터들의 정방향 표와 역방향 표,
반사판 표, 모두 아홉 장을 레지스터로 옮긴다. 한 클럭에 한 칸씩 234클럭이 든다(\ref{sec:e-loader}절).
릴레이식에서 배선 카드를 꽂는 일에 해당한다. 엔진은 한 걸음 더 나아가, 릴레이식이라면 위치마다 시프터가 새로
만들었을 스크램블러의 치환표를 미리 계산해 분산 RAM에 넣어 두고 읽어 쓴다. 다음 절의 주제다. 그 메모리는
메뉴의 선마다 하나씩 사본을 둔다(\ref{sec:e-ram}절). 릴레이식이 메뉴의 선마다 유닛을 하나씩 두는 것과 닮았다.

\section{자원과 시간을 맞바꾼다}\label{sec:tradeoff}

회로 설계에서 가장 자주 하는 결정은 ``부품을 더 쓰고 빨리 끝낼 것인가, 부품을 아끼고 오래 걸릴 것인가''다.
릴레이식에서는 거의 하지 않던 결정이다.

\Relay 릴레이식에는 이 거래의 여지가 거의 없다. 이어진 선을 전류로 찾으려면 메뉴의 연결이 \emph{모두 동시에}
있어야 한다(\ref{sec:f-current}절). 그래서 메뉴의 선마다 유닛이 하나씩 있어야 하고(\ref{sec:c-scrambler}절),
유닛 하나를 여러 선에 차례로 돌려 쓸 수 없다. 위치 하나의 시간도 릴레이의 개수가 아니라 릴레이가 넘어가고 붙는
시간이 정한다. 순서기의 한 주기, 63~ms다(\ref{sec:t-sequencer}절). 설계에서 저울질한 것은 릴레이 수와 시간이
아니라 릴레이 수와 다른 대가였다. 접점이 네 회로인 범용 릴레이를 쓰면 시프터의 릴레이가 절반(유닛당 260개)이
되지만, 코일 전력이 약 0.9~W여서 유닛 하나가 평균 100~W를 넘고 동작도 다섯 배 느리다(보고서의 릴레이 선정).
그래서 작은 릴레이를 많이 쓰는 쪽을 택했다. 뺄셈기 대신 차이를 직접 세는 계수기(\ref{sec:relay-counter}절)는
시간을 내주지 않고 구조만 바꿔 릴레이를 아낀 예다.

\Why 같은 계산을 하는 회로는 하나가 아니다. 계산에 필요한 연산을 전부 따로따로 부품으로 놓으면 한 클럭에
끝나지만 부품이 많이 든다. 부품 하나를 여러 클럭에 걸쳐 돌려 쓰면 부품은 적게 들지만 클럭이 여러 번
필요하다. 프로그램에서는 이 선택이 없다. 처리 장치가 하나이므로 언제나 뒤쪽이다.

FPGA에서 이 선택이 열리는 까닭은 \ref{sec:f-cost}절의 비교에 있다. 실물 기계의 위치 하나에는 연결을 한 번
바꾸고 전류를 한 번 거는 일밖에 없어서 쪼개어 나눠 쓸 시간이 없다. FPGA에서는 위치마다 클럭이 여러 번 들고,
클럭은 나노초 단위로 싸다. 그래서 회로 한 벌을 여러 클럭에 걸쳐 돌려 쓰고 부품을 아낄 수 있다.

\General 수 여덟 개의 합을 구한다고 하자(그림~\ref{fig:g-tradeoff}). 왼쪽은 덧셈기 일곱 개를 나무 모양으로
엮어 한 클럭에 답을 낸다. 오른쪽은 덧셈기 하나와 누산 레지스터 하나로 여덟 클럭에 걸쳐 하나씩 더한다.
신호 처리의 필터도 같다. 탭이 여덟 개인 필터는 곱셈기 여덟 개로 한 클럭에 계산할 수도 있고, 곱셈기 하나로
여덟 클럭에 걸쳐 계산할 수도 있다. 표본이 들어오는 속도가 느리면 뒤쪽으로 충분하다.

\begin{figure}[ht]
\centering
\begin{tikzpicture}[x=1cm,y=1cm,
   add/.style={draw,circle,fill=cbluebg,draw=cblue,minimum size=6mm,inner sep=0pt,font=\small}]
  \node[font=\small\bfseries] at (2.6,3.4) {병렬: 덧셈기 7개, 1클럭};
  \foreach \i in {0,...,7} { \node[font=\ttfamily\scriptsize] (x\i) at (\i*0.75,2.6) {x\i}; }
  \foreach \k/\a/\b in {0/0/1,1/2/3,2/4/5,3/6/7} {
    \node[add] (s\k) at (\k*1.5+0.375,1.7) {+};
    \draw[->] (x\a) -- (s\k); \draw[->] (x\b) -- (s\k);
  }
  \node[add] (t0) at (1.125,0.8) {+}; \node[add] (t1) at (4.125,0.8) {+};
  \draw[->] (s0) -- (t0); \draw[->] (s1) -- (t0); \draw[->] (s2) -- (t1); \draw[->] (s3) -- (t1);
  \node[add] (u) at (2.625,-0.1) {+};
  \draw[->] (t0) -- (u); \draw[->] (t1) -- (u);
  \draw[->] (u) -- ++(0,-0.6) node[below,font=\small] {합};
  \begin{scope}[xshift=9.0cm]
  \node[font=\small\bfseries] at (2.0,3.4) {직렬: 덧셈기 1개, 8클럭};
  \node[font=\ttfamily\small] (xin) at (-0.2,1.6) {x[i]};
  \node[add,minimum size=8mm] (p) at (1.3,1.6) {+};
  \node[reg,minimum width=14mm] (acc) at (3.3,1.6) {누산};
  \pic at (acc.south) {clkmark};
  \draw[->] (xin) -- (p); \draw[->] (p) -- (acc);
  \draw[->] (acc.east) -- ++(0.5,0) |- (2.3,0.5) -| (p.south);
  \draw[->] (acc.east) ++(0.5,0) -- ++(0.6,0) node[right,font=\small] {합};
  \node[note] at (2.0,-0.3) {클럭마다 수 하나를 더한다};
  \end{scope}
\end{tikzpicture}
\caption{같은 계산의 두 가지 구현.}
\label{fig:g-tradeoff}
\end{figure}

\Ours 봄베에서 이 선택은 스크램블러의 치환표를 언제, 어떻게 계산하는가에 나타난다.

실제 기계를 떠올려 보자. 봄베의 드럼 세 개 가운데 빠른 드럼은 계속 돌지만, 가운데 드럼은 빠른 드럼이 한 바퀴
돌 때 한 칸씩만, 느린 드럼은 그보다 더 드물게 움직인다. 26위치 동안 스크램블러의 ``안쪽''(가운데 로터, 왼쪽 로터, 반사판)은 그대로다.

\begin{itemize}
\item \textbf{기준 RTL}은 이 사실을 이용하지 않는다. 메뉴의 선마다 스크램블러 회로 전체(그림~\ref{fig:g-scr26})를
하나씩 두고 드럼 위치가 바뀔 때마다 26글자의 치환을 전부 다시 계산한다. 릴레이식과 같은 구조다. 릴레이식도
메뉴의 선마다 유닛을 하나씩 두고, 위치마다 시프터가 접점을 넘겨 치환을 새로 만든다. 릴레이식에서 그 값은
위치마다 접점이 넘어가는 시간이었고 그 시간은 어차피 63~ms 안에 들어 있었다. FPGA에서는 값이 공간으로 온다.
구조가 단순해서 맞는지 확인하기 쉽지만, 스크램블러 하나가 LUT 19,717개이므로 선이 16개인 메뉴에서는
스크램블러만으로 약 31만 5천 개다. 대상 칩(약 13만 4천 개)에 들어가지 않는다.
\item \textbf{FPGA 엔진}은 바뀌는 것만 다시 계산한다. 가운데 로터가 바뀔 때에만 안쪽 치환과 오른쪽 로터
26위치의 치환표를 여러 클럭에 걸쳐 계산해 분산 RAM에 넣어 두고, 위치마다 표에서 읽는다. 계산 회로는 한
벌만 있으면 되고 나머지는 메모리다.
\item \textbf{Rust 코어}도 같은 생각으로 짜여 있다. 가운데 로터가 바뀔 때만 표 26개를 다시 만든다
(\ref{ch:rust-bombe}장).
\end{itemize}

그림~\ref{fig:g-engine-idea}는 두 구조의 차이다. 엔진은 표를 만드는 데 클럭을 쓰지만 그것은 26위치에 한
번이다. 시간을 조금 내주고 LUT를 크게 돌려받았다. 엔진 전체는 LUT 41,040개에 들어간다. 엔진의 내부는
G부에서 다룬다.

\begin{figure}[ht]
\centering
\begin{tikzpicture}[x=1cm,y=1cm]
  \node[font=\small\bfseries] at (2.8,3.0) {기준 RTL: 위치마다 전부 다시 계산};
  \foreach \i/\yy in {0/2.1,1/1.3,2/0.5} {
    \node[blkb,minimum width=38mm,minimum height=6mm,font=\footnotesize] (s\i) at (2.2,\yy) {스크램블러 회로 (26글자분)};
    \draw[->] (s\i.east) -- ++(1.0,0);
  }
  \node at (2.2,-0.1) {$\vdots$};
  \node[note] at (2.2,-0.8) {메뉴의 선마다 한 벌 (최대 16벌)};
  \node[blkn,minimum width=10mm,minimum height=24mm] (sp) at (5.7,1.3) {전파};
  \begin{scope}[xshift=8.6cm]
  \node[font=\small\bfseries] at (3.1,3.0) {FPGA 엔진: 바뀔 때만 계산해 표에 저장};
  \node[blkb,minimum width=26mm,minimum height=10mm,font=\footnotesize] (calc) at (0.9,1.6) {치환 계산 회로\\(한 벌)};
  \node[reg,minimum width=22mm,minimum height=10mm,font=\footnotesize] (ram) at (3.9,1.6) {치환표\\(분산 RAM)};
  \node[blkn,minimum width=10mm,minimum height=24mm] (sp2) at (6.2,1.3) {전파};
  \draw[->] (calc) -- (ram); \draw[->] (ram) -- (sp2.west|-ram);
  \node[note] at (0.9,0.6) {가운데 로터가\\바뀔 때만 동작};
  \node[note] at (3.9,0.6) {위치마다 읽기};
  \end{scope}
\end{tikzpicture}
\caption{치환표를 얻는 두 가지 방법(개념도). 왼쪽은 릴레이식처럼 선마다 한 벌을 두고 위치마다 새로 만든다.}
\label{fig:g-engine-idea}
\end{figure}

같은 거래는 한 층 위에도 있다. 엔진 한 대는 로터 순서 하나를 맡으므로, 엔진을 두 대 두면 로터 순서 둘을
동시에 돈다. 배치와 배선을 마친 결과 한 대가 LUT 43,100개(32.2\%), 두 대가 83,257개(62.2\%)였고, 세 대이면 약
93\%가 되어 두 대로 정했다(\ref{sec:e-results}절). 릴레이식의 프레임에 유닛 자리가 12개뿐인 것처럼, 칩이라는
그릇의 크기가 나란히 둘 수 있는 수를 정한다.

\section{자원은 비싼 일에 든다}

두 기계에서 부품이 어디에 드는지 보면 \ref{sec:f-cost}절이 그대로 드러난다.

\begin{itemize}
\item \textbf{릴레이식}은 릴레이의 약 9할이 시프터, 곧 \emph{연결을 바꾸는} 회로다. 이어진 선을
찾는 데에는 릴레이가 하나도 들지 않는다. 메뉴 케이블과 대각선 보드는 구리이고, 릴레이가 드는 것은 찾은
결과를 읽고 판정하는 테스트 레지스터(감지 릴레이 26개와 판정·타이밍 릴레이 몇 개)뿐이다.
\item \textbf{FPGA 엔진}은 연결을 바꾸는 일을 표 읽기로 싸게 끝낸다. 크기를 좌우한 것은 \emph{이어진 선을
찾는} 전파 회로였다. 그 회로를 어떻게 적느냐에 따라 LUT가 약 40만 4천 개에서 41,040개로 달라졌다(\ref{ch:lesson}장).
\end{itemize}

물리가 공짜로 해 주던 일을 FPGA에서는 부품으로 짓는다. 그래서 FPGA 설계에서 자원을 아끼는 일은 대부분 그
일을 얼마나 싸게 짓는가의 문제다.

\section{정리}

\begin{itemize}
\item 릴레이식의 크기는 릴레이의 개수, 기판, 프레임, 코일의 열로 보인다. FPGA의 크기는 칩 안에 미리 깔린
부품을 몇 개 쓰는가다. 논리는 LUT에, 기억은 플립플롭에 담긴다. 회로의 크기는 주로 LUT의 수로 잰다.
\item LUT 하나는 입력 여섯 개까지의 함수다. 접점과 달리 식의 모양이 아니라 출력 하나가 의존하는 입력의 수가
크기를 정한다.
\item 변수로 무엇인가를 고르는 코드는 선택 회로이고, 고를 대상의 수에 비례해 커진다. 접점 나무가 $n$글자에
접점 $n-1$개를 쓰고 시프터가 릴레이식에서 가장 비쌌던 것과 같은 셈법이다. 큰 선택은 구조에 맞춰 작은 선택으로
나누면 싸진다.
\item 릴레이식의 표는 배선 카드와 드럼 기판의 구리여서 공짜이고 양방향이다. FPGA의 표는 메모리에 담기며,
되풀이되는 계산을 표로 만들어 두면 싸다.
\item 릴레이식은 연결이 모두 동시에 있어야 해서 자원과 시간을 맞바꿀 여지가 거의 없다. FPGA는 같은 계산을
부품을 많이 써서 빨리 할 수도, 적게 써서 천천히 할 수도 있다. 무엇이 자주 바뀌고 무엇이 드물게 바뀌는지를
보면 어디서 아낄지가 보인다.
\end{itemize}
